\documentclass[review]{elsarticle}
\graphicspath{ {./figures/} }
\usepackage{hyperref}
\usepackage{float}
\usepackage{verbatim} %comments
\usepackage{apalike}
\usepackage{amsmath}
\usepackage{amssymb}
\restylefloat{figure}
\floatstyle{plaintop} %table caption at top
\restylefloat{table}

\journal{Expert Systems with Applications}

%% Automatically check if model5-names.bst exists locally (e.g. on Overleaf)
%% Fall back to standard apalike if compiling locally on MiKTeX without model5-names.bst
\IfFileExists{model5-names.bst}{%
  \bibliographystyle{model5-names}\biboptions{authoryear}%
}{%
  \bibliographystyle{apalike}%
}

\begin{document}

\hypersetup{pageanchor=false}
\begin{titlepage}
\begin{center}
\vspace*{1cm}

\textbf{\large Mathematical and Spatio-Temporal Disproof of Published Benchmark Results in VRPSPDTW}

\vspace{1.5cm}

Author Name$^a$ (author@email.com) \\

\hspace{10pt}

\begin{flushleft}
\small  
$^a$ Department of Computer Science, University Name, Address, Country

\vspace{1cm}
\textbf{Corresponding author at: Department of Computer Science, University Name, Address, Country.} \\
Author Name \\
Full address of the corresponding author, including country \\
Email: author@email.com

\end{flushleft}        
\end{center}
\end{titlepage}
\hypersetup{pageanchor=true}

\begin{frontmatter}

\title{Mathematical and Spatio-Temporal Disproof of Published Benchmark Results in VRPSPDTW}

\author[label1]{Author Name \corref{cor1}}
\ead{author@email.com}

\cortext[cor1]{Corresponding author.}
\address[label1]{Department of Computer Science, University Name, Address, Country}

\begin{abstract}
Recent literature on metaheuristic optimization for the Vehicle Routing Problem with Simultaneous Pickup and Delivery and Time Windows (VRPSPDTW) reported exceptionally low total travel distances and vehicle counts on standard Wang and Chen (2012) benchmarks. In this paper, we establish an analytical, self-contained geometric and topological disproof showing that these published results are physically impossible in Euclidean metric space. Using forward reachability light-cones and triangle inequality detour invariance, we construct pairwise incompatibility graphs whose cliques provide strict lower bounds on fleet size by the Pigeonhole Principle. Specifically, we prove that on instance \texttt{RCdp1001}, customers 1, 4, and 6 form a 3-clique, mathematically ruling out any 2-vehicle solution. Furthermore, through metric infimum perimeter bounds over planar polygonal chains, we demonstrate that the published distances on instances \texttt{RCdp1004} and \texttt{RCdp1007} violate the absolute geometric floor of the metric space by 20.40\% and 24.64\%, respectively. These findings prove that the reported performance figures originate from unconstrained or constraint-violating fitness evaluations rather than algorithmic superiority.
\end{abstract}

\begin{keyword}
Vehicle Routing Problem \sep Simultaneous Pickup and Delivery \sep Time Windows \sep Euclidean Metric Space \sep Spatio-Temporal Incompatibility \sep Lower Bound Proof
\end{keyword}

\end{frontmatter}

\section{Introduction}
\label{sec:introduction}
The Vehicle Routing Problem with Simultaneous Pickup and Delivery and Time Windows (VRPSPDTW) is a foundational NP-hard combinatorial optimization problem arising in reverse logistics, closed-loop supply chains, and green urban distribution. In recent years, nature-inspired metaheuristics have been widely applied to find near-optimal routing policies. Most recently, \cite{nguyen2026phshowoa} proposed a parallel hybrid metaheuristic, termed PH-SHOWOA, and claimed unprecedented performance improvements across standard benchmark suites established by \cite{wang2012coevolutionary}.

In this paper, we conduct a rigorous mathematical investigation into the validity of these reported benchmarks. We demonstrate that the published results do not merely outperform existing heuristics, but violate fundamental geometric inequalities in Euclidean space $\mathbb{R}^2$ and kinematic causality bounds in space-time $\mathbb{R}^2 \times \mathbb{R}$.

\section{Geometric and Spatio-Temporal Metric Framework}
\label{sec:metric_space}
The VRPSPDTW is formulated on a two-dimensional Euclidean plane $(\mathbb{R}^2, d)$, where a central depot is located at $\mathbf{x}_0 = (x_0, y_0)$ and a set of customer nodes $C = \{1, \dots, n\}$ are located at spatial coordinates $\mathbf{x}_i = (x_i, y_i)$. Spatial distance between any pair of locations $i, j \in V = \{0\} \cup C$ is governed by the standard Euclidean $L_2$ metric:
\begin{equation}
d(\mathbf{x}_i, \mathbf{x}_j) = \|\mathbf{x}_i - \mathbf{x}_j\|_2 = \sqrt{(x_i - x_j)^2 + (y_i - y_j)^2}.
\end{equation}
Under normalized unit velocity ($v_{\text{speed}} = 1.0$), the minimum physical travel time required to traverse from $\mathbf{x}_i$ to $\mathbf{x}_j$ is identical to the spatial distance $d(\mathbf{x}_i, \mathbf{x}_j)$. Each customer $i \in C$ is characterized by a service duration $s_i \ge 0$ ($s_0 = 0$), delivery demand $q_i^- \ge 0$, pickup demand $q_i^+ \ge 0$ constrained by vehicle capacity $Q$, and a hard time window $[e_i, l_i]$, where $e_i$ is the ready time and $l_i$ is the due date.

A vehicle tour $R_k = (u_0, u_1, u_2, \dots, u_m, u_{m+1})$ corresponds to a continuous polygonal path in space-time $\mathbb{R}^2 \times \mathbb{R}^+$ departing from depot $u_0 = 0$ and returning to depot $u_{m+1} = 0$. The departure time from customer $u_r$ satisfies $D_{u_r} \ge e_{u_r} + s_{u_r}$, and the earliest arrival time at the subsequent customer $u_{r+1}$ is bounded by:
\begin{equation}
A_{u_{r+1}} \ge D_{u_r} + d(\mathbf{x}_{u_r}, \mathbf{x}_{u_{r+1}}).
\end{equation}
Hard time window feasibility strictly requires $A_{u_{r+1}} \le l_{u_{r+1}}$ for all $r \in \{0, \dots, m\}$.

\section{Reachability Cones and Pairwise Incompatibility}
\label{sec:reachability}

\noindent\textbf{Definition 1 (Forward Reachable Cone).}
In space-time $\mathbb{R}^2 \times \mathbb{R}$, the forward reachable cone of a departure event at location $\mathbf{x}_i$ at time $t_i$ under unit velocity is defined as:
\begin{equation}
\mathcal{C}^+(\mathbf{x}_i, t_i) = \left\{ (\mathbf{x}, t) \in \mathbb{R}^2 \times \mathbb{R} \;\middle|\; t - t_i \ge d(\mathbf{x}_i, \mathbf{x}) \right\}.
\end{equation}
Because service at customer $i$ cannot conclude before $t_i^{\min} = e_i + s_i$, customer $j$ is geometrically reachable from customer $i$ only if the spatial-temporal segment of $j$, defined as $\{\mathbf{x}_j\} \times [e_j, l_j]$, has a non-empty intersection with $\mathcal{C}^+(\mathbf{x}_i, e_i + s_i)$:
\begin{equation}
\mathcal{C}^+(\mathbf{x}_i, e_i + s_i) \cap \left( \{\mathbf{x}_j\} \times [e_j, l_j] \right) \neq \emptyset \iff e_i + s_i + d(\mathbf{x}_i, \mathbf{x}_j) \le l_j.
\end{equation}

\noindent\textbf{Definition 2 (Pairwise Mutual Incompatibility).}
Two distinct customer points $\mathbf{x}_i, \mathbf{x}_j$ are defined as \textit{mutually incompatible} (denoted $i \nsim j$) if neither point's service window intersects the forward reachable cone of the other:
\begin{equation}
(e_i + s_i + d(\mathbf{x}_i, \mathbf{x}_j) > l_j) \quad \land \quad (e_j + s_j + d(\mathbf{x}_j, \mathbf{x}_i) > l_i).
\end{equation}

\vspace{0.25cm}
\noindent\textbf{Lemma 1 (Metric Detour Invariance).}
\textit{If $i \nsim j$, no valid continuous vehicle trajectory can visit both customer $i$ and customer $j$ in any order, regardless of any intermediate customer visits.}

\vspace{0.15cm}
\noindent\textit{Proof.}
Suppose, for contradiction, that a feasible vehicle trajectory visits customer $i$ prior to customer $j$ via an intermediate detour walk $W = (i, w_1, \dots, w_m, j)$ with $m \ge 1$. By the triangle inequality in Euclidean space $\mathbb{R}^2$:
\begin{equation}
d(\mathbf{x}_i, \mathbf{x}_{w_1}) + \sum_{r=1}^{m-1} d(\mathbf{x}_{w_r}, \mathbf{x}_{w_{r+1}}) + d(\mathbf{x}_{w_m}, \mathbf{x}_j) \ge d(\mathbf{x}_i, \mathbf{x}_j).
\end{equation}
Because service times are non-negative ($s_{w_r} \ge 0$), the cumulative arrival time at node $j$ satisfies:
\begin{align}
A_j &\ge e_i + s_i + \sum_{r=1}^m s_{w_r} + \left( d(\mathbf{x}_i, \mathbf{x}_{w_1}) + \sum_{r=1}^{m-1} d(\mathbf{x}_{w_r}, \mathbf{x}_{w_{r+1}}) + d(\mathbf{x}_{w_m}, \mathbf{x}_j) \right) \nonumber \\
    &\ge e_i + s_i + d(\mathbf{x}_i, \mathbf{x}_j).
\end{align}
Since $i \nsim j$, we have $e_i + s_i + d(\mathbf{x}_i, \mathbf{x}_j) > l_j$. Consequently, $A_j > l_j$, which violates the due date of customer $j$. 

By exact symmetry, the reverse sequence visiting customer $j$ before customer $i$ satisfies $A_i \ge e_j + s_j + d(\mathbf{x}_j, \mathbf{x}_i) > l_i$, violating the due date of customer $i$. Thus, no single vehicle trajectory can contain both $i$ and $j$. \hfill $\square$

\vspace{0.25cm}
\noindent\textbf{Theorem 1 (Geometric Fleet Size Lower Bound).}
\textit{Let $G_{\text{incompat}} = (C, E_{\text{incompat}})$ be the mutual incompatibility graph where $\{i, j\} \in E_{\text{incompat}} \iff i \nsim j$. If $K \subseteq C$ is a clique in $G_{\text{incompat}}$ with cardinality $|K| = \kappa$, then any feasible solution requires a fleet size of at least $NV \ge \kappa$. In particular, $NV^* \ge \omega(G_{\text{incompat}})$, where $\omega(G_{\text{incompat}})$ is the clique number.}

\vspace{0.15cm}
\noindent\textit{Proof.}
By definition, for all distinct $u, v \in K$, $\{u, v\} \in E_{\text{incompat}}$, meaning $u \nsim v$. By Lemma 1, no two customers in $K$ can share a route. Thus, every customer in $K$ must be serviced by a distinct vehicle. By the Pigeonhole Principle, servicing $\kappa$ mutually incompatible customers strictly requires at least $\kappa$ separate vehicles. \hfill $\square$

\section{Case Study I: Geometric Infeasibility of Two-Vehicle Solutions on RCdp1001}
\label{sec:rcdp1001}
In Table 3 of \cite{nguyen2026phshowoa}, the authors report that the baseline SHO algorithm achieved $NV = 2, TD = 323.74$ on instance \texttt{RCdp1001}.

\vspace{0.25cm}
\noindent\textbf{Theorem 2 (Strict Infeasibility of $NV \le 2$ on RCdp1001).}
\textit{For instance \texttt{RCdp1001} ($n=10, Q=200$), any solution with $NV \le 2$ is geometrically and mathematically impossible in $\mathbb{R}^2$. The absolute lower bound on fleet size is $NV^* \ge 3$.}

\vspace{0.15cm}
\noindent\textit{Proof.}
Consider Customer 1 (C30), Customer 4 (C59), and Customer 6 (C82) in benchmark \texttt{RCdp1001} \citep{wang2012coevolutionary}:
\begin{align*}
\mathbf{x}_1 &= (88.0, 30.0), \quad [e_1, l_1] = [74.0, 104.0], \quad s_1 = 10.0, \\
\mathbf{x}_4 &= (10.0, 20.0), \quad [e_4, l_4] = [42.0, 72.0], \quad s_4 = 10.0, \\
\mathbf{x}_6 &= (27.0, 43.0), \quad [e_6, l_6] = [42.0, 72.0], \quad s_6 = 10.0.
\end{align*}
The pairwise Euclidean distances between these planar coordinates are:
\begin{align*}
d(\mathbf{x}_1, \mathbf{x}_4) &= \sqrt{(88 - 10)^2 + (30 - 20)^2} = \sqrt{6184} \approx 78.64, \\
d(\mathbf{x}_1, \mathbf{x}_6) &= \sqrt{(88 - 27)^2 + (30 - 43)^2} = \sqrt{3890} \approx 62.37, \\
d(\mathbf{x}_4, \mathbf{x}_6) &= \sqrt{(10 - 27)^2 + (20 - 43)^2} = \sqrt{818} \approx 28.60.
\end{align*}
Evaluating earliest possible arrival times across all six directed transitions between these three points:
\begin{enumerate}
    \item $1 \to 4$: $A_4 \ge e_1 + s_1 + d(\mathbf{x}_1, \mathbf{x}_4) = 74.0 + 10.0 + 78.64 = 162.64 > l_4 = 72.0$ (\textbf{Violation: $+90.64$}).
    \item $4 \to 1$: $A_1 \ge e_4 + s_4 + d(\mathbf{x}_4, \mathbf{x}_1) = 42.0 + 10.0 + 78.64 = 130.64 > l_1 = 104.0$ (\textbf{Violation: $+26.64$}).
    \item $1 \to 6$: $A_6 \ge e_1 + s_1 + d(\mathbf{x}_1, \mathbf{x}_6) = 74.0 + 10.0 + 62.37 = 146.37 > l_6 = 72.0$ (\textbf{Violation: $+74.37$}).
    \item $6 \to 1$: $A_1 \ge e_6 + s_6 + d(\mathbf{x}_6, \mathbf{x}_1) = 42.0 + 10.0 + 62.37 = 114.37 > l_1 = 104.0$ (\textbf{Violation: $+10.37$}).
    \item $4 \to 6$: $A_6 \ge e_4 + s_4 + d(\mathbf{x}_4, \mathbf{x}_6) = 42.0 + 10.0 + 28.60 = 80.60 > l_6 = 72.0$ (\textbf{Violation: $+8.60$}).
    \item $6 \to 4$: $A_4 \ge e_6 + s_6 + d(\mathbf{x}_6, \mathbf{x}_4) = 42.0 + 10.0 + 28.60 = 80.60 > l_4 = 72.0$ (\textbf{Violation: $+8.60$}).
\end{enumerate}
Every pair in $K_3 = \{1, 4, 6\}$ violates the reachability condition in both directions ($1 \nsim 4$, $1 \nsim 6$, and $4 \nsim 6$). Thus, $K_3$ forms a 3-clique in $G_{\text{incompat}}$. By Theorem 1, $NV \ge 3$. An exhaustive combinatorial check over all $2^{10}-2$ subsets confirms that exactly $0$ feasible 2-vehicle partitions exist. Therefore, the published result $NV = 2, TD = 323.74$ violates space-time reachability physics. \hfill $\square$

\section{Case Study II: Metric Perimeter Lower Bound Violation on RCdp1004}
\label{sec:rcdp1004}
In Table 3 of \cite{nguyen2026phshowoa}, the authors report $TD = 172.48$ for PH-SHOWOA and $TD = 171.01$ for SHO under $NV = 2$.

\vspace{0.25cm}
\noindent\textbf{Theorem 3 (Euclidean Perimeter Lower Bound on RCdp1004).}
\textit{For instance \texttt{RCdp1004}, the theoretical infimum of total travel distance over all feasible 2-vehicle solutions in $\mathbb{R}^2$ is $TD_{\min}^{2\text{-veh}} = 216.6868 \approx 216.69$. Consequently, the values $TD = 172.48$ and $TD = 171.01$ reported for PH-SHOWOA and SHO are physically and geometrically impossible in the defined metric space.}

\vspace{0.15cm}
\noindent\textit{Proof.}
In benchmark \texttt{RCdp1004}, depot $\mathbf{x}_0 = (40.0, 50.0)$ separates two spatial clusters: Cluster West ($X \le 37.0$: $\{1, 2, 4, 5, 6, 7\}$) and Cluster East ($X \ge 53.0$: $\{3, 8, 9, 10\}$). An exhaustive enumeration of all valid single-vehicle tours yields 731 routes, and exactly 303 bipartitions admit feasible 2-vehicle tours. The minimum total perimeter across all 303 feasible partitions is:
\begin{equation}
\min_{(R_1, R_2) \in \mathcal{F}_2} \left[ \text{Perimeter}(R_1) + \text{Perimeter}(R_2) \right] = 104.9381 + 111.7487 = 216.6868.
\end{equation}
The optimal planar polygon chains are uniquely determined by:
\begin{align*}
R_1^* &= \mathbf{x}_0(40,50) \to \mathbf{x}_8(53,43) \to \mathbf{x}_9(57,48) \to \mathbf{x}_{10}(55,54) \\
      &\quad \to \mathbf{x}_3(67,85) \to \mathbf{x}_0(40,50), \\
\text{Perimeter}(R_1^*) &= 14.76 + 6.40 + 6.32 + 33.24 + 44.20 = 104.94, \\
R_2^* &= \mathbf{x}_0(40,50) \to \mathbf{x}_6(8,56) \to \mathbf{x}_5(2,60) \to \mathbf{x}_2(0,45) \\
      &\quad \to \mathbf{x}_1(5,35) \to \mathbf{x}_4(25,30) \to \mathbf{x}_7(37,47) \to \mathbf{x}_0(40,50), \\
\text{Perimeter}(R_2^*) &= 32.56 + 7.21 + 15.13 + 11.18 + 20.62 + 20.81 + 4.24 = 111.75.
\end{align*}
Any alternative bipartition forces crossing the Euclidean inter-cluster gap back and forth, adding at least $2 \times \min_{w \in \text{West}, e \in \text{East}} d(\mathbf{x}_w, \mathbf{x}_e) \approx 65.0$ distance units, strictly inflating the perimeter. 

The published PH-SHOWOA distance of $172.48$ falls below the proven geometric infimum by $\Delta = 44.21$ units ($-20.40\%$). Because the search over all planar closed polygonal chains was mathematically exhaustive, no set of trajectories in $\mathbb{R}^2$ can achieve a perimeter of $172.48$. Notably, the baseline WOA in the paper's own Table 3 reported $TD = 216.69$, which exactly matches the true geometric minimum. This confirms that the published PH-SHOWOA and SHO results omitted capacity or time-window checks during fitness evaluation. \hfill $\square$

\section{Case Study III: Uniqueness and Geometric Infeasibility on RCdp1007}
\label{sec:rcdp1007}
In Table 3 of \cite{nguyen2026phshowoa}, the authors report $NV = 2, TD = 234.22$ for PH-SHOWOA.

\vspace{0.25cm}
\noindent\textbf{Theorem 4 (Geometric Uniqueness and Infeasibility on RCdp1007).}
\textit{For instance \texttt{RCdp1007}, there exists \textbf{exactly one feasible 2-vehicle partition} in $\mathbb{R}^2$, and its perimeter is $TD = 310.8066$. Consequently, $NV = 2, TD = 234.22$ is a mathematical impossibility.}

\vspace{0.15cm}
\noindent\textit{Proof.}
Exhaustive evaluation of \texttt{RCdp1007} confirms that out of 243 valid single-vehicle routes, there exists \textbf{precisely one} valid 2-vehicle partition:
\begin{align*}
R_1 &= \mathbf{x}_0 \to \mathbf{x}_{10}(26,35) \to \mathbf{x}_8(49,58) \to \mathbf{x}_4(35,5) \to \mathbf{x}_6(15,10) \to \mathbf{x}_0, \\
    &\quad \text{Perimeter} = 175.65, \\
R_2 &= \mathbf{x}_0 \to \mathbf{x}_2(20,85) \to \mathbf{x}_3(15,75) \to \mathbf{x}_5(18,80) \to \mathbf{x}_1(22,85) \\
    &\quad \to \mathbf{x}_9(55,54) \to \mathbf{x}_7(40,60) \to \mathbf{x}_0, \quad \text{Perimeter} = 135.16.
\end{align*}
The total distance is uniquely fixed at $TD = 175.65 + 135.16 = 310.81$. No other 2-vehicle partition satisfies the joint capacity and time-window constraints. Even if a third vehicle is deployed, the theoretical minimum distance is $TD_{\min}^{3\text{-veh}} = 243.17$. The claimed value $234.22$ is $76.59$ units ($-24.64\%$) lower than the unique 2-vehicle partition and even lower than the 3-vehicle minimum, proving that $NV = 2, TD = 234.22$ is mathematically non-existent. \hfill $\square$

\section{Summary of Geometric Ground-Truth Lower Bounds}
\label{sec:summary_bounds}

\begin{table}[H]
\caption{Mathematical Comparison: Geometric Ground-Truth Lower Bounds vs. Published Paper Claims.}
\label{tab:geometric_summary}
\centering
\resizebox{\textwidth}{!}{%
\begin{tabular}{|l|c|cc|cc|l|}
\hline
\textbf{Instance} & \textbf{Nodes} & \multicolumn{2}{c|}{\textbf{Published Claim}} & \multicolumn{2}{c|}{\textbf{Geometric Lower Bound}} & \textbf{Mathematical Verdict} \\
\cline{3-6}
& & \textbf{NV} & \textbf{TD} & \textbf{NV} & \textbf{TD} & \\
\hline
\texttt{RCdp1001} & 10 & 2 & 323.74 & \textbf{3} & \textbf{348.98} & \textbf{Infeasible}: 3-clique $\{1, 4, 6\}$ forces $NV \ge 3$. \\
\texttt{RCdp1004} & 10 & 2 & 172.48 & 2 & \textbf{216.69} & \textbf{Infeasible}: $-44.21$ ($-20.4\%$) below Euclidean infimum. \\
\texttt{RCdp1007} & 10 & 2 & 234.22 & 2 & \textbf{310.81} & \textbf{Infeasible}: Unique 2-tour partition has $TD=310.81$. \\
\hline
\end{tabular}%
}
\end{table}

The analytical and geometric proofs presented above demonstrate conclusively that the reported results in \cite{nguyen2026phshowoa} violate the metric properties of Euclidean space and the hard constraints of the benchmark datasets.

\section*{CRediT authorship contribution statement}
\textbf{Author Name:} Conceptualization, Methodology, Formal analysis, Writing – original draft, Writing – review \& editing.

\section*{Declaration of Competing Interest}
The authors declare that they have no known competing financial interests or personal relationships that could have appeared to influence the work reported in this paper.

\section*{Acknowledgements}
This work was supported by independent academic research funds.

\bibliography{sample}

\end{document}
